\documentclass[11pt,a4paper]{article}

\usepackage[a4paper,margin=1in]{geometry}
\usepackage{times}
\usepackage{hyperref}
\usepackage{graphicx}
\usepackage{titlesec}
\usepackage{enumitem}
\usepackage{amsmath}

\setlist{noitemsep}
\setlength{\parskip}{0.6em}

\title{\textbf{FindBestCareers: An LLM-Powered Multi-Agent System for Intelligent Career Development and Job Matching}}

\author{
    HE Yufeng (UID: 3036661094)
    \and
    LIU Jinlin (UID: 3036656245)
}

\date{COMP7607 | Natural Language Processing | Project Report}

\begin{document}

\maketitle

\section{Introduction and Motivation}

With the rapid development of large language models (LLMs) and their increasing adoption in industry, there is a growing interest in applying such models to real-world, end-to-end workflows. One particularly impactful area is the recruitment pipeline, where both job seekers and employers are overwhelmed by the volume and heterogeneity of information: long-form resumes in PDF format, noisy job descriptions, and unstructured interview interactions.

Traditional recruitment systems mainly rely on keyword matching or heuristic ranking, which struggle to capture nuanced skill profiles, cross-domain competencies, and candidate--job fit at a deeper semantic level. At the same time, modern LLMs exhibit strong capabilities in natural language understanding, summarization, and dialogue, making them ideal candidates for supporting resume analysis, job recommendation, and interview simulation. However, deploying such systems in a practical setting requires careful system design, robust data processing, and alignment between frontend UI and backend APIs.

In this project, we design and implement \textbf{FindBestCareers}, an LLM-powered web system that integrates three major functions into a unified pipeline:

\begin{itemize}
    \item \textbf{Resume parsing and skill profiling}: parsing PDF resumes, extracting structured information, and scoring skills using a combination of LLM prompting and a curated skill taxonomy.
    \item \textbf{Job matching}: enriching job postings with skill tags and computing a match score between a candidate's profile and available positions.
    \item \textbf{AI interview simulation}: conducting a structured, three-phase interview (self-introduction, technical Q\&A, and feedback) with robust state tracking and LLM-based evaluation.
\end{itemize}

The system is built as a full-stack application. The backend is implemented in Python using Flask, Pandas, PyPDF2, and an abstraction layer over multiple LLM providers (OpenAI and DeepSeek). The frontend is implemented in React (Vite + TypeScript) with a modern UI that exposes the three main functionalities. Throughout the project, we focus on \emph{faithful} use of LLMs (grounded in an explicit skill taxonomy), robust error handling in API calls, and a clear separation between infrastructural components (e.g., LLM client) and application logic.

\section{System Overview}

The overall system architecture is organized into three cooperating subsystems: (1) resume understanding and skill tagging, (2) job enrichment and matching, and (3) interview orchestration and dialogue. A thin Flask API server mediates all communication between the frontend and backend logic.

\subsection{Backend Components}

The backend is implemented in Python and built around the following core modules:

\begin{itemize}
    \item \textbf{LLM client abstraction} (\texttt{llm\_client.py}): a unified client that can route chat requests to different providers (OpenAI, DeepSeek) and models, based on a configuration file and/or environment variables. It also integrates a key rotation mechanism via an \texttt{APIKeyManager} for robustness.
    \item \textbf{Resume parser} (\texttt{resume\_parser.py}): responsible for PDF extraction, skill taxonomy loading, LLM prompting for category selection, and final skill scoring.
    \item \textbf{Job matcher} (\texttt{job\_matcher.py}): parses job skill tags from a pre-enriched CSV and computes a job--candidate match score.
    \item \textbf{Interview agent} (\texttt{interview\_agent.py}): implements a three-phase interview state machine with structured prompts and JSON parsing for question generation and answer evaluation.
    \item \textbf{API server} (\texttt{api\_server.py}): a Flask-based server exposing RESTful endpoints for resume upload, job querying and matching, and interview session management.
\end{itemize}

The backend adopts an in-memory store for resumes and interview sessions for simplicity, which is sufficient for a course project and local demonstration. In a production setting, this layer can be replaced by a persistent database.

\subsection{Frontend Components}

The frontend is implemented with React, TypeScript, and Tailwind CSS, bundled using Vite. It is organized into three main pages:

\begin{itemize}
    \item \textbf{ResumePage}: handles PDF file upload, displays extracted basic information, and visualizes skill scores grouped by categories.
    \item \textbf{JobsPage}: lists job positions, shows a match score for each job, provides skill-based filtering and pagination, and offers an ``Apply'' button that opens the original job URL.
    \item \textbf{InterviewPage}: manages the AI interview session, displaying a chat-style interface between the user and the AI interviewer across the three phases.
\end{itemize}

Navigation is implemented with a custom \texttt{Navigation} component that switches between these pages while keeping their internal state mounted (via \texttt{hidden} classes instead of unmounting), ensuring that context (e.g., filters, page indices, conversation history) is preserved when users switch between views.

\section{Resume Parsing and Skill Profiling}

\subsection{PDF Ingestion and Basic Extraction}

Upon receiving a PDF file upload from the frontend, the backend (\texttt{/api/resume/upload}) saves the file to a designated \texttt{uploads/} directory and invokes the \texttt{ResumeParser} module. We use PyPDF2 to extract raw text from each page, followed by lightweight cleaning (e.g., removing repeated headers/footers and normalizing whitespace).

In addition to the raw text, we maintain a structured container for high-level resume fields, including:

\begin{itemize}
    \item Personal information: name, email, phone.
    \item Education history: school, degree, major, graduation year.
    \item Work experience: company, position, duration, and free-form description.
\end{itemize}

These fields are inferred via LLM prompting. The LLM is asked to output a constrained JSON object describing this structure, which is then parsed and stored. Robust parsing is implemented to handle the common case where the LLM returns JSON wrapped in markdown code blocks or additional commentary.

\subsection{Skill Taxonomy and Tag Extraction}

A key design decision in this project is to ground skill evaluation in a pre-defined taxonomy rather than letting the LLM hallucinate arbitrary skills. The taxonomy is stored in a CSV file (\texttt{all\_labels.csv}) and contains, for each tag:

\begin{itemize}
    \item A \texttt{level\_3rd} category (e.g., ``Machine Learning Engineer'').
    \item One or more concrete tags (e.g., ``Python'', ``PyTorch'', ``data preprocessing'').
    \item Metadata indicating default or core tags.
\end{itemize}

The resume parser performs two LLM-based steps:

\begin{enumerate}
    \item \textbf{Category selection}: The LLM is given the full set of unique \texttt{level\_3rd} values and the resume text, and is asked to select up to 10 categories that best describe the candidate's background and interests. This step allows the model to integrate information from both the resume and, later, the candidate's self-description in the interview.
    \item \textbf{Tag scoring}: Given the selected categories, the corresponding tag sets, plus default tags and simple string matches in the resume text, we build a candidate pool of skill tags. The LLM is then asked to assign each tag a score from 1 to 5, following detailed scoring guidelines (e.g., evidence of real projects, explicit metrics, years of experience).
\end{enumerate}

To keep the output consistent with the frontend, we post-process the LLM's responses into a list of \texttt{\{skill\_name, score\}} objects. Skills are further categorized into buckets such as ``Core Skills'', ``Language / Tools'', and ``Other'', and sorted so that core skills (e.g., high scores or manually marked as core) appear at the top. This ordering is directly reflected in the frontend skill visualization.

\subsection{Frontend Visualization and UX Considerations}

\begin{itemize}
    \item \texttt{/api/jobs}: to load all jobs and their parsed required skills.
    \item \texttt{/api/jobs/match}: to obtain match scores for the current resume (if available).
\end{itemize}

As shown in Fig~\ref{fig:resume}, the frontend \texttt{ResumePage} displays both the high-level resume fields and the scored skills. For each skill, we show:
\begin{figure}
    \centering
    \includegraphics[width=0.9\linewidth]{resume.png}
    \caption{The rendering of ResumePage.}
    \label{fig:resume}
\end{figure}
\begin{itemize}
    \item A textual name (e.g., ``PyTorch'').
    \item A 1--5 star visualization of the score.
    \item Grouping by category, with the ``Other'' category intentionally placed at the bottom to reduce visual clutter.
\end{itemize}

From a UX perspective, we also ensure that:
\begin{itemize}
    \item Once a resume is uploaded and parsed, the structured result is kept in memory, so users can navigate to other pages (e.g., job matching, interview) and come back without losing their analysis.
    \item The resume ID is stored in \texttt{localStorage} under \texttt{current\_resume\_id} so that other modules (job matching, interview) can reference the same candidate profile.
\end{itemize}

\section{Job Enrichment and Matching}
The rendering of this part can be found in Fig.~\ref{fig:position}.
\begin{figure}
    \centering
    \includegraphics[width=0.9\linewidth]{position.png}
    \caption{The rendering of Job Enrichment and Matching.}
    \label{fig:position}
\end{figure}
\subsection{Job Data and Skill Tags}

Job postings are loaded from a pre-enriched CSV file (\texttt{bytedance\_jobs\_enriched.csv}). Each row corresponds to one position and contains fields such as company name, job title, location, description, and a compact \texttt{skill\_tags} string. The latter encodes a sequence of skills and scores in a custom format (e.g., ``C++, 5 ; NLP , 5'').

On the backend, the API endpoint \texttt{/api/jobs} uses a helper function \texttt{parse\_skill\_tags} to convert the \texttt{skill\_tags} string into a simple list of skill names. Additional metadata---such as job family, minimum degree requirement, and application URL---are also extracted and exposed to the frontend.

\subsection{Matching Algorithm}

The matcher (\texttt{job\_matcher.py}) computes candidate-job alignment via weighted skill overlap. For each job, we parse \texttt{skill\_tags} into \((s, w)\) pairs (skill, importance weight), compute overlapping skills with the candidate profile (case-insensitive with normalization), and aggregate a match score as Eq.~\ref{eq:score}:

\begin{equation}
    \text{Score}(c, j) = \alpha \cdot |\mathcal{S}_c \cap \mathcal{S}_j| + (1-\alpha) \cdot \frac{1}{|\mathcal{S}_c \cap \mathcal{S}_j| \cdot maxscore} \sum_{s \in \mathcal{S}_c \cap \mathcal{S}_j} score(s), \label{eq:score}
\end{equation}

where \(\mathcal{S}_c\) and \(\mathcal{S}_j\) are candidate and job skill sets, \(maxscore\) is the max socre of every skill, in this work, this number is setted as 5, \(score(s)\) is candidate's job score on skill \(s\), and \(\alpha=0.6\) prioritizes coverage. Scores normalize to 0--100 and rank via \texttt{job\_matcher.py}.


\subsection{Frontend Job Browser and Skill Filtering}

The page supports several UX features:

\begin{itemize}
    \item \textbf{Skill-based filtering}: Instead of using a fixed, hard-coded skill list, we automatically extract all unique \texttt{required\_skills} from the loaded jobs and display them as filter keywords. Users can select multiple skills; jobs must satisfy all selected keywords (conjunctive filter).
    \item \textbf{Keyword search}: Users can search by job title, company, or description.
    \item \textbf{Pagination}: Jobs are displayed in pages with a small number of items per page for readability. When the total number of pages is large, the paginator shows only the first few page numbers by default, with an option to expand and reveal all page numbers.
    \item \textbf{Compact filters}: Since the number of unique skill keywords can be large, the UI shows only the first 50 by default, with an ``expand/collapse'' toggle to reveal or hide the full list.
\end{itemize}

Each job card shows the job title, company, location, a textual description, the computed match score (with a color-coded badge), and a list of required skills. An ``Apply'' button uses the job's \texttt{apply\_url} field to open the original application page in a new browser tab. 

\section{AI Interview Module}

\begin{figure}
    \centering
    \includegraphics[width=0.9\linewidth]{interview.png}
    \caption{The rendering of AI Interview Module.}
    \label{fig:interview}
\end{figure}
\subsection{Three-Phase Interview Design}

As shown in Fig.~\ref{fig:interview}, the AI interview module is designed around a deterministic three-phase state machine:

\begin{enumerate}
    \item \textbf{Self-introduction (Greeting)}: The system greets the candidate in English and invites them to introduce their background, experience, and target position. No technical question is asked in this phase.
    \item \textbf{Technical Q\&A}: A fixed number of technical questions (five by default) are asked sequentially. For each question, the system evaluates the candidate's answer and then generates the next question.
    \item \textbf{Summary and feedback}: After the last technical question, the system generates a comprehensive feedback report summarizing the candidate's performance, strengths, weaknesses, and suggested next steps.
\end{enumerate}

The interview agent keeps track of the current phase, the number of Q\&A rounds completed, and the entire conversation history (questions, answers, and evaluations). This state is stored in an in-memory session object.

\subsection{Prompting and JSON Parsing}

To improve the robustness of LLM responses, the \texttt{InterviewAgent} uses structured prompts and explicitly requests strict JSON output for machine-readable parts of the flow:

\begin{itemize}
    \item \textbf{Technical question generation}: The system prompt describes the desired characteristics of a question (e.g., specific, medium difficulty, requiring project-level experience). The LLM is asked to output a JSON object \texttt{\{"question": ..., "difficulty": ..., "category": ...\}} which is parsed and turned into an \texttt{InterviewQuestion} dataclass instance.
    \item \textbf{Answer evaluation}: Similarly, the LLM is asked to return a JSON object with fields such as score, feedback, strengths, and improvements. This is parsed into an \texttt{AnswerEvaluation} dataclass.
\end{itemize}

However, LLMs occasionally wrap JSON in markdown code blocks or add natural language commentary. To handle this, we implement a helper \texttt{\_extract\_json\_from\_text} that first looks for JSON within fenced code blocks, and if that fails, searches for the first JSON object in the text. If parsing still fails, we fall back to reasonable default messages and questions.

A unified \texttt{LLMClient} is used to send these prompts, so switching between OpenAI and DeepSeek (or changing model names) requires only configuration changes rather than code modifications in the interview logic.

\subsection{Frontend Chat Interface}

The \texttt{InterviewPage} presents a chat-style interface that closely resembles common messaging applications. Key design aspects include:

\begin{itemize}
    \item \textbf{State persistence}: The interview session and chat history are kept alive even when the user navigates away to other pages and back, thanks to the persistent mounting of components in \texttt{App.tsx}.
    \item \textbf{Typing indicator}: When the backend is processing a response, the UI displays an animated typing indicator to provide feedback and reduce perceived latency.
    \item \textbf{Input handling}: Pressing Enter sends the message; Shift+Enter inserts a newline. The send button is disabled while waiting for the AI response or when the input is empty.
    \item \textbf{Session termination}: A ``End Interview'' button resets the frontend state and stops further messages from being sent.
\end{itemize}

From the user's perspective, the AI interviewer first thanks them for their introduction, then proceeds with numbered technical questions, each followed by a short evaluation and a transition to the next question, concluding with a summarized feedback report.

\section{Implementation Details}

\subsection{LLM Provider Abstraction}

To decouple the application logic from LLMs, we introduce a unified client in \texttt{llm\_client.py}. Configuration is controlled via a JSON file (\texttt{llm\_config.json}):

\begin{itemize}
    \item \texttt{provider}: \texttt{"openai"} or \texttt{"deepseek"}.
    \item \texttt{model}: model identifier, e.g., \texttt{"gpt-4o-mini"} or \texttt{"deepseek-chat"}.
    \item \texttt{timeout}, \texttt{max\_retry}, \texttt{temperature}: inference parameters.
\end{itemize}

This design allows the same codebase to be tested with different models in a flexible way and aligns with the assignment's requirement to support both OpenAI and DeepSeek.

\section{Evaluation and Discussion}

\subsection{Functional Testing}

Due to time and resource constraints, we focus on qualitative evaluation and functional correctness:

\begin{itemize}
    \item \textbf{Resume parsing}: We tested the system with several English and Chinese resumes. Basic fields (name, education, experience) were extracted reliably. Skill scores broadly aligned with our expectations, especially for clearly stated technical skills.
    \item \textbf{Job matching}: The match scores assigned higher values to jobs whose skill tags overlapped more with the candidate's core skills, and the sorted ranking appeared meaningful. For example, a candidate with strong C++ and client development skills was ranked higher for C++ client intern roles.
    \item \textbf{AI interview}: The three-phase structure worked as designed: greeting and explicit explanation of stages, followed by five distinct technical questions, and then a final feedback report. We observed that enforcing JSON output significantly reduced prompt-following errors and hallucinated formatting.
\end{itemize}

\subsection{Lessons Learned}

From an engineering perspective, several design decisions proved beneficial:

\begin{itemize}
    \item Treating the LLM as a \emph{component} with clear contracts (input prompts and output JSON), rather than as a monolithic black box, greatly simplified debugging and made the system more robust.
    \item Separating low-level LLM infrastructure from application logic allowed us to swap models and providers without rewriting business code.
    \item Iterative refinement based on real error logs (e.g., handling malformed JSON, duplicate questions) was essential; many issues only surfaced during interactive testing.
\end{itemize}

Overall, the project achieves the main goals set at the beginning of the semester: integrating LLMs into a full-stack application, addressing concrete engineering challenges, and reflecting on the opportunities and limitations of LLM-based systems in realistic scenarios.

\end{document}